% Look!  An Appendix!

\section{Review of Differential Topology}

In this appendix, we try to provide the reader a short introduction to differential topology, sufficient to read the thesis. 

\subsection{Vectors, Covectors and Vector Bundles}

In order to define the notion of tangent spaces, and cotangent spaces, we first define the notion of function germs. This definition captures the behavior of a function $f$ around a point $p$. The following presentation is from \textcite{Bjorn2018}

\begin{defn}[A Function Germ]

Let $M$, $N$ be smooth manifolds, and consider $p \in M$. We say that functions $f: M \to N$ and $g: M \to N$ are equivalent $f \sim g$, if there exists an open neighborhood of $p, U_{fg}$ such that $f = g$ on $U_{fg}$. In that case, we denote the equivalence class of $f$, $[f] = \overline{f}$, and the function germ $\overline{f}$ is:

$$\overline{f}: (M, p) \to (N, f(p))$$
\end{defn}

\noindent Now, the tangent space to $M$ at $p$ is the equivalence class of function germs passing through $p$ with a given velocity:

\begin{defn}[Tangent space to $M$ at $p$, $T_{p} M$]
Let $M$ be a smooth $n$-dimensional manifold. Let $p \in M$. 

$$W_{p} = \{\operatorname{germs } \; \overline{\gamma}: (\mathbb{R}, 0) \to (M, p) \}$$

\noindent Two germs $\overline{\gamma}, \overline{\gamma}_{1} \in W_{p}$ are equivalent if two function germs $\overline{\phi}: (M, p) \to (\mathbb{R}, \phi(p))$ we have that $(\phi \gamma)'(0) = (\phi \gamma_{1})'(0)$. We define the tangent space to $M$ at $p$ to be the set of equivalence classes 

$$T_{p} M = W_{p}/\sim $$
\end{defn}

\begin{remark}
Therefore, the tangent space is best thought of as an equivalence class of curves
\end{remark}

As above, let $M$ be a smooth manifold, and consider $p \in M$. We can define the cotangent space $T_{p}^{*} M$ as the vector space dual to $T_{p} M$. 

\begin{defn}[Dual Vector Space]
Given a real vector space $V$, we define $V^{*}$ to be the vector space $\operatorname{Hom}_{\mathbb{R}}(V, \mathbb{R})$, the set of $\mathbb{R}$-linear maps from $V \to \mathbb{R}$
\end{defn}

\begin{defn}[Cotangent Space]
As above, let $M$ be a smooth manifold, and consider $p \in M$. Given the tangent vector space $T_{p} M$, the cotangent space, $T^{*}_{p} M$ contains all linear maps $p$ of the form:
$$T_{p}^{*} M  = \{p: V \to \mathbb{R}, p \text{ linear } \}$$
\end{defn}

\noindent Now, given the tangent spaces and cotangent spaces, there are two set-theoretic constructions we often make: $T^{*}M = \bigsqcup_{m \in M} T^{*}_{m} M$ and $TM = \bigsqcup_{m \in M} T_{m} M$. These are sets which are disjoint unions of all tangent planes as we index $m \in M$. These sets are known as the cotangent and tangent bundles, respectively. 

\begin{defn}[A Real Vector Bundle]
A real vector bundle of rank \( n \) is a pair \( (E, X) \), where \( E \) is a topological space space called the total space, and \( X \) is a topological space called the base space, together with a continuous map \( \pi : E \to X \) such that:

\begin{enumerate}
    \item For each \( p \in X \), the fibre \( \pi^{-1}(p) \) is a \( n \)-dimensional real vector space.

    \item For every point \( p \in X \), there exists an open neighborhood \( U \subset X \) containing \( p \), and a homeomorphism
    \[
    h : \pi^{-1}(U) \to U \times \mathbb{R}^n
    \]
    such that the following diagram commutes:
    \[
    \begin{tikzcd}
    \pi^{-1}(U) \arrow[rr, "h"] \arrow[dr, "\pi|_{\pi^{-1}(U)}"'] & & U \times \mathbb{R}^n \arrow[dl, "\text{proj}_1"] \\
    & U &
    \end{tikzcd}
    \]
    In other words, \( \pi = \text{proj}_1 \circ h \).

    \item For each \( q \in U \), the map \( h \) is to a vector space isomorphism on fibres:
    \[
    h_q : \pi^{-1}(q) \to \{q\} \times \mathbb{R}^n \xrightarrow{\cong} \mathbb{R}^n,
    \]
\end{enumerate}
\end{defn}

\begin{remark}
A real vector bundle is a fibreing over the base space by vector spaces, such that we can move through the fibreing continuously via $h$.  
\end{remark}

\begin{remark}
$T^{*} M$ and $TM$ are both vector bundles
\end{remark}

\noindent \textbf{References}: \textcite{Terek2020}, \textcite{Bjorn2018}

\subsection{Tensor Product}
Given two vector real spaces $V$ and $W$, the tensor product of $V$ and $W$ is:
$$V \otimes W = (V \times W)/S$$
\noindent where $S$ is the linear subspace generated by linear combinations of
\begin{align}
e_{1} &= (v_{1} + v_{2}, w) - (v_{1}, w) - (v_{2}, w) \\ 
e_{2} &= (v, w_{1} + w_{2})  - (v, w_{1}) - (v, w_{2}) \\ 
e_{3} &= (sv, w) - s(v, w) \\
e_{4} &= (v, sw) - s(v, w)
\end{align}

\begin{defn}[Tensor Product of Linear Maps]
The tensor product of two linear maps $A: V \to V$ and $B: W \to W$ is the unique linear map such that:

$$(A \otimes B): (V \otimes W) \to (V \otimes W)$$

\noindent where $s \in \mathbb{R}$, $(v,w) \in V \times W$. 
\end{defn}

\noindent \textbf{References}: \textcite{Isham:1999rh}

\subsection{Exterior Derivative, Closed and Exact Forms}

\begin{defn}
\noindent An $n$-form is a tensor field $\omega$ of type $(0, n)$ that is totally skew-symmetric, in the sense that, for any permutation $P$ of the indices $1, 2, \dots, n$:

\[
\omega(X_{1}, \dots, X_{n}) = (-1)^{\deg P} \omega(X_{P(1)}, \dots, X_{P(n)})
\]
\end{defn}

\noindent In particular, the $\wedge$ product is the tool which defines multiplication of $n_{1}$-forms by $n_{2}$-forms and outputs $(n_{1} + n_{2})$-forms. It is the antisymmetric part of the tensor product of both forms:

\[
\omega_{1} \wedge \omega_{2} = \frac{1}{n_{1}! n_{2}!} \sum_{\sigma \in P} (-1)^{\deg \sigma} (\omega_{1} \otimes \omega_{2})^{\sigma}
\]

\begin{defn}[Exterior Derivative $d$]
The exterior derivative of $\omega$ is the $(n+1)$-form $d \omega$ defined by:

\begin{align}
d\omega(X_{1}, \dots, X_{n+1}) = \sum_{i=1}^{n+1} (-1)^{i+1} X_{i}(\omega(X_{1}, \dots, \hat{X}_{i}, \dots, X_{n+1}))  \\
+ \sum_{i < j} (-1)^{i+j} \omega([X_{i}, X_{j}], X_{1}, \dots, \hat{X}_{i}, \dots, \hat{X}_{j}, \dots, X_{n+1})
\end{align}
\end{defn} 

\begin{remark}
\noindent In the case of $2$-forms, we have that:

\[
d \omega(X,Y) = X(\langle \omega, Y \rangle) - Y(\langle \omega, X \rangle) - \langle \omega, [X,Y] \rangle
\]
\end{remark}

\noindent In the DeRham Theory, we consider the DeRham complex, a sequence of vector spaces:

\[
0 \xrightarrow{d} C^{\infty}(\mathcal{M}) \xrightarrow{d} A^{1}(\mathcal{M}) \xrightarrow{d} A^2(\mathcal{M}) \xrightarrow{d} \dots \xrightarrow{d} A^{m}(\mathcal{M}) \xrightarrow{d} 0
\]

\noindent Notably, since $d^2 = 0$, one has:

\[
\text{Im}(d: A^{n-1} \to A^{n}) \subset \text{Ker}(d: A^{n} \to A^{n+1})
\]

\noindent An $n$-form $\omega$ is closed if $d \omega = 0$. We define the set of all closed $n$-forms as:

\[
Z^{n}(\mathcal{M}) = \text{Ker}(d: A^{n} \to A^{n+1}) = \{\omega \in A^{n}(\mathcal{M}) \; | \; d\omega = 0\}
\]

\noindent Similarly, the set of all exact forms is:

\[
B^{n} = \text{Im}(d: A^{n-1} \to A^{n}) = \{\omega \in A^{n}(\mathcal{M}) \; | \; \omega = d \beta \text{ for some } \beta \in A^{n-1}(\mathcal{M}) \}
\]

\[
H_{DR}^{n}(\mathcal{M}) = Z^{n}(\mathcal{M})/B^{n}(\mathcal{M})
\]

\noindent \textbf{References}: \textcite{Isham:1999rh}

\section{Additional Topics}

\subsection{Derivation of A Lie Bracket Identity}

\begin{defn}[Lie Derivative of a Vector Field]
For a smooth manifold $M$, consider $X, Y \in \mathfrak{X}(M)$. Then, the Lie derivative of $X$, in the $Y$ direction is:

$$\mathcal{L}_{Y}(X) = [Y,X]$$

\noindent Where $[X, Y]$ denotes the vector field commutator of $X$ and $Y$.  
\end{defn}

\noindent For the following results, we apply Cartan's Magic Formula:

\begin{theorem}[Cartan's Magic Formula]

$$\mathcal{L}_{X} = d(i_{X}) + i_{X} d$$

\end{theorem}

\begin{lemma}[Fundamental Contraction Identity]

Let $\mathcal{L}$ be the Lie-derivative, $i$ denote contraction, and $d$ denote the exterior derivative. Then:

\label{FundamentalContracIdentity}
$$[\mathcal{L}_{X}, i_{Y}] = i_{[X, Y]}$$
\end{lemma}

\begin{proof}
Expand $\mathcal{L}_{X}$ according to the Lie-bracket: 

$$[\mathcal{L}_{X}, i_{Y}] = [d \circ i_{X}  + i_{X} \circ d, i_{Y}] = [i_{X} \circ d, i_{Y}] = i_{X} \mathcal{L}$$

\noindent As $[i_{X}, i_{Y}] = 0$. Then we use the fact that:

$$i_{x} \circ d i_{Y} - i_{Y} \circ i_{X} \circ d = i_{x} \circ d i_{Y} - i_{Y} \circ i_{X} \circ d = i_{X}( d i_{Y} + i_{Y} d) = i_{X} \mathcal{L}_{Y}$$
\end{proof}

Now using this lemma, we prove the following:

\begin{lemma}
\label{PoissonIsLie}
$$i_{[X_{f}, X_{g}]} \omega = \mathcal{L}_{X_{f}} i_{X_{g}} \omega $$
\end{lemma}

\begin{proof}
Using \ref{FundamentalContracIdentity} we have that:

$$i_{[X_{f}, X_{g}]} \omega = [\mathcal{L}_{X_{f}}, i_{X_{g}}]$$

\noindent Then:

$$i_{[X_{f}, X_{g}]} \omega = \mathcal{L}_{X_{f}}(dg) = d(X_{f}(g)) = -\omega(X_{f}, X_{g})$$

\noindent Therefore, we conclude.
\end{proof}

\noindent \textbf{References}: \textcite{Silva2001}, \textcite{Isham:1999rh}

\subsection{Example with Fubini Study Metric on $\mathbb{C} P^{1}$}

\begin{example}[Fubini-Study Metric on $\mathbb{C}P^1$]
The Fubini-Study metric is the area form on a sphere $S^{2}$ when $V = \mathbb{C}^2$, and hence $\mathbb{P}V = \mathbb{C}P^{1}$.
For $\mathbb{C}P^1$, we use homogeneous coordinates $[z_0:z_1]$. In the chart where $z_0 \neq 0$, we can use the coordinate $w = z_1/z_0$. A point in $\mathbb{C}P^1$ is represented as $v = (1, w)$ in this chart, with inner product: $\langle v, v \rangle = 1 + |w|^2$. Let's denote our two tangent vectors as: $\zeta_1 = (0, \xi_1)$, $\zeta_2 = (0, \xi_2)$. Now, $\langle \zeta_1, \zeta_2 \rangle = \overline{\xi_1} \cdot \xi_2 = \overline{\xi_1} \xi_2$ and $\langle \zeta_1, v \rangle = \overline{\xi_1} \cdot w = \overline{\xi_1} w$ and $\langle v, \zeta_2 \rangle = \overline{w} \cdot \xi_2 = \overline{w} \xi_2$. Then, the form is:
\begin{equation}
\omega_{[v]}(\tilde{\zeta}_1, \tilde{\zeta}_2) = \text{Im} \left( \frac{(\overline{\xi_1} \xi_2)(1 + |w|^2) - (\overline{\xi_1} w)(\overline{w} \xi_2)}{(1 + |w|^2)^2} \right)
\end{equation}
Note that the following hold:
\begin{align}
(\overline{\xi_1} \xi_2)(1 + |w|^2) - (\overline{\xi_1} w)(\overline{w} \xi_2) &= \overline{\xi_1} \xi_2 + \overline{\xi_1} \xi_2 |w|^2 - \overline{\xi_1} w \overline{w} \xi_2\\
&= \overline{\xi_1} \xi_2 + \overline{\xi_1} \xi_2 |w|^2 - \overline{\xi_1} \xi_2 |w|^2\\
&= \overline{\xi_1} \xi_2
\end{align}
Note that $w \overline{w} = |w|^2$, which is why $\overline{\xi_1} w \overline{w} \xi_2 = \overline{\xi_1} \xi_2 |w|^2$.
\begin{equation}
\omega_{[v]}(\tilde{\zeta}_1, \tilde{\zeta}_2) = \text{Im} \left( \frac{\overline{\xi_1} \xi_2}{(1 + |w|^2)^2} \right)
\end{equation}
\begin{equation}
\text{Im}(\overline{\xi_1} \xi_2) = \frac{1}{2i}(\overline{\xi_1} \xi_2 - \xi_1 \overline{\xi_2})
\end{equation}
Note that the expression for vectors $\overline{\xi_1} \xi_2 - \xi_1 \overline{\xi_2}$ implies that the symplectic form must be antisymmetric and contain a factor $- dw \wedge d\overline{w}$. This can be written in terms of differential forms as:
\begin{equation}
\omega_{[v]} = \frac{i/2}{(1 + |w|^2)^2} dw \wedge d\overline{w}
\end{equation}
Writing $w = x + iy$, we have that: $\omega_{FS} = \frac{1}{(1 + x^2 + y^2)^2} dx \wedge dy$. Compare this to the expression for the area form of $S^{2}$ in $\mathbb{R}^2$ using stereographic projection: $\omega = \frac{4}{(1 + x^2 + y^2)^2} dx \wedge dy$. Therefore, $\omega = \frac{1}{4} \omega_{FS}$ and up to a choice of constant in the definition, both agree. See here: \textcite{Silva2001}. 
\end{example}

\noindent \textbf{References}: \textcite{Silva2001}



\newpage
\setcounter{section}{5}








